% Part: first-order-logic
% Chapter: syntax-and-semantics
% Section: satisfaction

\documentclass[../../../include/open-logic-section]{subfiles}

\begin{document}

\olfileid{fol}{syn}{sat}

\olsection{Vervulling van \article{formula} \printtoken{S}{formula}
  in \article{structure} \printtoken{S}{structure}}

\begin{explain}
De grondbegrippen die enerzijds uitdrukkingen zoals termen en
!!{formula}s en anderzijds !!{structure}s met elkaar verbinden, zijn
de \emph{!!{value}} van een term en de \emph{vervulling} van
!!a{formula}. Informeel is de !!{value} van een term een !!{element} van
!!a{structure}---als de term enkel een constante is, is zijn !!{value}
het object dat de !!{structure} aan de constante toekent; als de term
met !!{function}s is opgebouwd, wordt de !!{value} berekend uit de
!!{value}s van constanten en de functies die aan de functies in de term
zijn toegekend. !!^a{formula} wordt in !!a{structure} \emph{vervuld} als
de interpretatie van de predicaten de !!{formula} waar maakt in het
domein van de !!{structure}. Dit vervullingsbegrip wordt inductief
vastgelegd: de specificatie van de !!{structure} bepaalt rechtstreeks
wanneer atomaire !!{formula}s worden vervuld, en voor een samengestelde
!!{formula} hangt dit af van de hoofdoperator en van de vraag of de
onmiddellijke !!{subformula}s worden vervuld.

Bij kwantoren ligt dit iets ingewikkelder. De onmiddellijke
!!{subformula} van een gekwantificeerde !!{formula} heeft namelijk een
vrije !!{variable}, terwijl !!{structure}s geen !!{value}s voor
!!{variable}s vastleggen. Om dit probleem op te lossen voeren we ook
\emph{bedelingen} in. We definiëren vervulling dan niet alleen ten
opzichte van !!a{structure}, maar ten opzichte van !!a{structure}
samen met !!a{variable} bedeling.
\end{explain}

\begin{defn}[Bedeling]
Een \emph{bedeling}~$s$ voor !!a{structure}~$\Struct{M}$ is een
functie die elke !!{variable} afbeeldt op !!a{element} van~$\Domain M$,
dat wil zeggen: $s\colon \Var \to \Domain M$.
\end{defn}

\begin{explain}
!!^a{structure} kent aan elke !!{constant} een !!{value} toe, en een
bedeling aan elke variabele. We willen echter dat ook de daarmee
opgebouwde termen !!{element}s van het !!{domain} aanduiden. Daarom
definiëren we de !!{value} van termen inductief. Voor !!{constant}s en
variabelen is de waarde precies wat de !!{structure} respectievelijk de
bedeling voorschrijft; bij samengestelde termen wordt zij recursief
berekend met de functies die de !!{structure} aan de !!{function}s
toekent.
\end{explain}

\begin{defn}[!!^{value} van termen]
Zij $t$ een term van de taal~$\Lang L$, zij $\Struct M$ een
!!{structure} voor~$\Lang L$, en zij $s$ !!a{variable} bedeling
voor~$\Struct M$. De \emph{!!{value}}~$\Value{t}{M}[s]$ wordt als
volgt gedefinieerd:
\begin{enumerate}
\item \indcase{t}{c}{$\Value{\indfrm}{M}[s] = \Assign{\indcomplex}{M}$.}
\item \indcase{t}{x}{$\Value{\indfrm}{M}[s] = s(\indcomplex)$.}
\item \indcase{t}{\Atom{f}{t_1, \ldots, t_n}}{
\[
\Value{\indfrm}{M}[s] = \Assign{f}{M}(\Value{t_1}{M}[s], \ldots,
\Value{t_n}{M}[s]).
\]}
\end{enumerate}
\end{defn}

\begin{defn}[$x$-variant]
Als $s$ !!a{variable} bedeling is voor !!a{structure}~$\Struct M$, dan heet
elke !!{variable} bedeling~$s'$ voor~$\Struct M$ die van~$s$ hoogstens
verschilt in wat zij aan~$x$ toekent, een \emph{$x$-variant} van~$s$. Als
$s'$ een $x$-variant van~$s$ is, schrijven we $\varAssign{s'}{s}{x}$.
\end{defn}

\begin{explain}
Merk op dat een $x$-variant van een bedeling~$s$ \emph{niet} noodzakelijk iets
anders aan~$x$ toekent. Elke bedeling geldt zelfs als een $x$-variant
van zichzelf.
\end{explain}

\begin{defn}
  Als $s$ !!a{variable} bedeling is voor !!a{structure}~$\Struct M$
  en $m \in \Domain{M}$, dan is de bedeling~$\Subst{s}{m}{x}$ de
  variabeletoekenning die wordt gedefinieerd door
  \[\Subst{s}{m}{x}(y) = \begin{cases}
    m & \text{als } y \ident x\\
    s(y) & \text{anders}.
  \end{cases}\]
\end{defn}

Met andere woorden: $\Subst{s}{m}{x}$ is de bijzondere $x$-variant van~$s$
die het !!{element}~$m$ van het domein aan~$x$ toekent en aan andere
!!{variable}s dan~$x$ hetzelfde toekent als $s$.

\begin{defn}[Vervulling]
\ollabel{defn:satisfaction}
!!^a{formula}~$!A$ wordt in !!a{structure}~$\Struct M$ onder
!!a{variable} bedeling~$s$ vervuld, in symbolen:
$\Sat{M}{!A}[s]$, volgens de volgende recursieve definitie. (We schrijven
$\Sat/{M}{!A}[s]$ voor ``niet $\Sat{M}{!A}[s]$.'')
\begin{enumerate}
\tagitem{prvFalse}{%
  \indcase{!A}{\lfalse}{$\Sat/{M}{\indfrm}[s]$.}}{}

\tagitem{prvTrue}{%
  \indcase{!A}{\ltrue}{$\Sat{M}{\indfrm}[s]$.}}{}

\item \indcase{!A}{\Atom{R}{t_1, \dots, t_n}}{$\Sat{M}{\indfrm}[s]$
  dan en slechts dan als $\langle \Value{t_1}{M}[s], \dots, \Value{t_n}{M}[s] \rangle \in
  \Assign{R}{M}$.}

\item \indcase{!A}{\eq[t_1][t_2]}{$\Sat{M}{\indfrm}[s]$ dan en slechts
  dan als $\Value{t_1}{M}[s] = \Value{t_2}{M}[s]$.}

\tagitem{prvNot}{%
  \indcase{!A}{\lnot !B}{$\Sat{M}{\indfrm}[s]$ dan en slechts dan als
    $\Sat/{M}{!B}[s]$.}}{}

\tagitem{prvAnd}{%
  \indcase{!A}{(!B \land !C)}{$\Sat{M}{\indfrm}[s]$ dan en slechts dan als $\Sat{M}{!B}[s]$
    en $\Sat{M}{!C}[s]$.}}{}

\tagitem{prvOr}{%
  \indcase{!A}{(!B \lor !C)}{$\Sat{M}{\indfrm}[s]$ dan en slechts dan als
    $\Sat{M}{!B}[s]$ of $\Sat{M}{!C}[s]$ (of beide).}}{}

\tagitem{prvIf}{%
  \indcase{!A}{(!B \lif !C)}{$\Sat{M}{\indfrm}[s]$ dan en slechts dan als $\Sat/{M}{!B}[s]$
    of $\Sat{M}{!C}[s]$ (of beide).}}{}

\tagitem{prvIff}{%
  \indcase{!A}{(!B \liff !C)}{$\Sat{M}{\indfrm}[s]$ dan en slechts dan als ofwel zowel
    $\Sat{M}{!B}[s]$ als $\Sat{M}{!C}[s]$, ofwel noch $\Sat{M}{!B}[s]$
    noch $\Sat{M}{!C}[s]$.}}{}

\tagitem{prvAll}{%
  \indcase{!A}{\lforall[x][!B]}{$\Sat{M}{\indfrm}[s]$ dan en slechts dan als voor elk
    !!{element}~$m \in \Domain M$, $\Sat{M}{!B}[\Subst{s}{m}{x}]$.}}{}

\tagitem{prvEx}{%
  \indcase{!A}{\lexists[x][!B]}{$\Sat{M}{\indfrm}[s]$ dan en slechts dan als voor ten minste
  één !!{element}~$m \in \Domain M$, $\Sat{M}{!B}[\Subst{s}{m}{x}]$.}}{}
\end{enumerate}
\end{defn}

\begin{explain}
De bedelingen zijn van belang in de laatste
\iftag{notprvEx,notprvAll}{clausule}{twee clausules}.\iftag{prvAll}{ We
kunnen de vervulling van $\lforall[x][!B(x)]$ niet definiëren met ``voor alle $m \in
\Domain{M}$, $\Sat{M}{!B(m)}$.''}{}\iftag{prvEx}{ We kunnen de
vervulling van $\lexists[x][!B(x)]$ niet definiëren met ``voor ten minste één $m \in
\Domain{M}$, $\Sat{M}{!B(m)}$.''}{} De reden is dat een $m \in
\Domain M$ geen symbool van de taal is. Daarom is $!B(m)$ geen
!!a{formula} (dat wil zeggen: $\Subst{!B}{m}{x}$ is ongedefinieerd). We
kunnen evenmin aannemen dat er !!{constant}s of termen beschikbaar zijn
die elk !!{element} van~$\Struct{M}$ benoemen, want de definitie van
!!{structure}s eist dit nergens. In de standaardtaal is de verzameling
!!{constant}s !!{denumerable}; als $\Domain{M}$ niet !!{enumerable} is,
zijn er dus niet eens genoeg !!{constant}s om elk object te benoemen.

We lossen dit probleem op met !!{variable} bedelingen, die variabelen
rechtstreeks met !!{element}s van het domein verbinden. In plaats van
bijvoorbeeld te zeggen dat
\iftag{prvEx}{$\lexists[x][!B(x)]$}{$\lforall[x][!B(x)]$} in~$\Struct M$
wordt vervuld dan en slechts dan als \iftag{prvEx}{voor ten minste één $m \in
\Domain{M}$ aan de zojuist bedoelde voorwaarde wordt voldaan}{voor alle $m \in \Domain{M}$, $\Sat{M}{!B(m)}$}, zeggen we dat de formule in~$\Struct M$
\emph{onder}~$s$ wordt vervuld dan en slechts dan als $!B(x)$ onder
$\Subst{s}{m}{x}$ wordt vervuld \iftag{prvEx}{voor ten minste
één}{voor elke} $m \in \Domain M$.
\end{explain}

\begin{ex}
Zij $\Lang{L} = \{a, b, f, R\}$, waarbij $a$ en $b$ !!{constant}s zijn,
$f$~een tweeplaatsige !!{function} is en $R$~een tweeplaatsig !!{predicate}.
Beschouw de !!{structure}~$\Struct{M}$ die als volgt is gedefinieerd:
\begin{enumerate}
\item $\Domain M = \{1, 2, 3, 4\}$
\item $\Assign{a}{M} = 1$
\item $\Assign{b}{M} = 2$
\item $\Assign{f}{M}(x, y) = x+y$ als $x+y \le 3$ en anders $= 3$.
\item $\Assign{R}{M} = \{\tuple{1, 1}, \tuple{1, 2}, \tuple{2, 3}, \tuple{2, 4}\}$
\end{enumerate}
De functie $s(x) = 1$, die aan elke !!{variable} het element
$1 \in \Domain{M}$ toekent, is een bedeling voor~$\Struct{M}$.

Dan geldt
\begin{align*}
\Value{f(a,b)}{M}[s] & = \Assign{f}{M}(\Value{a}{M}[s], \Value{b}{M}[s]).
\intertext{Omdat $a$ en $b$ !!{constant}s zijn, geldt $\Value{a}{M}[s]
  = \Assign{a}{M} = 1$ en $\Value{b}{M}[s] = \Assign{b}{M} = 2$. Dus}
\Value{f(a,b)}{M}[s] & = \Assign{f}{M}(1, 2) = 1+2 = 3.
\intertext{Om de waarde van $f(f(a,b),a)$ te berekenen, beschouwen we}
\Value{f(f(a,b),a)}{M}[s] & = \Assign{f}{M}(\Value{f(a, b)}{M}[s],
\Value{a}{M}[s]) = \Assign{f}{M}(3, 1) = 3,
\intertext{omdat $3+1 > 3$. Omdat $s(x) = 1$ en $\Value{x}{M}[s] =
  s(x)$, geldt ook}
\Value{f(f(a,b),x)}{M}[s] & = \Assign{f}{M}(\Value{f(a, b)}{M}[s],
\Value{x}{M}[s]) = \Assign{f}{M}(3, 1) = 3,
\end{align*}

Een atomaire !!{formula}~$R(t_1, t_2)$ wordt vervuld als het tupel van
de waarden van haar argumenten, dus $\tuple{\Value{t_1}{M}[s],
  \Value{t_2}{M}[s]}$, !!a{element} is van~$\Assign{R}{M}$. Zo geldt
$\Sat{M}{R(b,f(a,b))}[s]$, omdat $\tuple{\Value{b}{M},
  \Value{f(a,b)}{M}} = \tuple{2, 3} \in \Assign{R}{M}$, maar
$\Sat/{M}{R(x, f(a,b))}[s]$, omdat $\tuple{1, 3} \notin \Assign{R}{M}$.

Om vast te stellen of een niet-atomaire formule~$!A$ wordt vervuld, pas
je de clausule van de inductieve definitie toe die bij haar
hoofdoperator hoort. De hoofdoperator van $R(a, a) \lif (R(b,
x) \lor R(x, b))$ is bijvoorbeeld~$\lif$, en
\begin{align*}
  & \Sat{M}{R(a, a) \lif (R(b, x) \lor R(x, b))}[s] \text{ dan en slechts dan als }\\
  & \qquad
  \Sat/{M}{R(a,a)}[s] \text{ of } \Sat{M}{R(b, x) \lor R(x, b)}[s]
  \intertext{Omdat $\Sat{M}{R(a,a)}[s]$ (want $\tuple{1,1} \in
    \Assign{R}{M}$), kunnen we het antwoord nog niet bepalen. Eerst
    moeten we nagaan of $\Sat{M}{R(b, x) \lor R(x, b)}[s]$:}
  & \Sat{M}{R(b, x) \lor R(x, b)}[s] \text{ dan en slechts dan als }\\
  & \qquad \Sat{M}{R(b,x)}[s] \text{ of } \Sat{M}{R(x, b)}[s]
  \intertext{Dit is inderdaad het geval, want $\Sat{M}{R(x, b)}[s]$
    (omdat $\tuple{1,2} \in \Assign{R}{M}$).}
\end{align*}

Een $x$-variant van~$s$ is, zoals gezegd, een bedeling die van $s$
hoogstens verschilt in wat zij aan~$x$ toekent. Voor elk !!{element} van
~$\Domain{M}$ bestaat er een $x$-variant van~$s$:
\begin{align*}
  s_1 & = \Subst{s}{1}{x}, &
  s_2 & = \Subst{s}{2}{x},\\
  s_3 & = \Subst{s}{3}{x}, &
  s_4 & = \Subst{s}{4}{x}.
\end{align*}
Zo is bijvoorbeeld $s_2(x) = 2$ en $s_2(y) = s(y) = 1$ voor alle variabelen~$y$
behalve~$x$. Dit zijn alle $x$-varianten van~$s$ voor de
structuur~$\Struct{M}$, want $\Domain{M} = \{1, 2, 3, 4\}$. In het
bijzonder geldt $s_1 = s$ ($s$~is altijd een $x$-variant van zichzelf).

\iftag{prvEx}{Om vast te stellen of een existentieel gekwantificeerde
  !!{formula}~$\lexists[x][!A(x)]$ wordt vervuld, moeten we nagaan
  of $\Sat{M}{!A(x)}[\Subst{s}{m}{x}]$ voor ten minste één $m \in \Domain
  M$. Dus
  \[
  \Sat{M}{\lexists[x][(R(b,x) \lor R(x,b))]}[s],
  \]
  want $\Sat{M}{R(b,x) \lor R(x, b)}[\Subst{s}{1}{x}]$
  ($\Subst{s}{3}{x}$ zou eveneens volstaan). Daarentegen geldt
  \[
  \Sat/{M}{\lexists[x][(R(b,x) \land R(x,b))]}[s]
  \]
  want welke $m \in \Domain{M}$ we ook kiezen, $\Sat/{M}{R(b,x) \land R(x,b)}[\Subst{s}{m}{x}]$.}{}

\iftag{prvAll}{Om vast te stellen of een universeel gekwantificeerde
  !!{formula}~$\lforall[x][!A(x)]$ wordt vervuld, moeten we nagaan
  of $\Sat{M}{!A(x)}[\Subst{s}{m}{x}]$ voor alle $m \in \Domain M$. Dus
  \[
  \Sat{M}{\lforall[x][(R(x,a) \lif R(a,x))]}[s],
  \]
  want $\Sat{M}{R(x,a) \lif R(a,x)}[\Subst{s}{m}{x}]$ voor alle $m \in
  \Domain M$. Voor $m = 1$ geldt $\Sat{M}{R(a,x)}[\Subst{s}{1}{x}]$,
  zodat het consequens waar is; voor $m = 2$, $3$ en~$4$ geldt
  $\Sat/{M}{R(x,a)}[\Subst{s}{m}{x}]$, zodat het antecedens onwaar is.
  Maar
  \[
  \Sat/{M}{\lforall[x][(R(a,x) \lif R(x,a))]}[s]
  \]
  want $\Sat/{M}{R(a,x) \lif R(x,a)}[\Subst{s}{2}{x}]$ (omdat
  $\Sat{M}{R(a, x)}[\Subst{s}{2}{x}]$ en $\Sat/{M}{R(x,
  a)}[\Subst{s}{2}{x}]$).}{}

\iftag{defEx}{Om vast te stellen of een existentieel gekwantificeerde
  !!{formula}~$\lexists[x][!A(x)]$ wordt vervuld, moeten we nagaan
  of $\Sat{M}{\lnot \lforall[x][\lnot !A(x)]}[s]$. Zo geldt
  \[
  \Sat{M}{\lexists[x][(R(b,x) \lor R(x,b))]}[s].
  \]
  Ten eerste geldt $\Sat{M}{R(b,x) \lor R(x, b)}[\Subst{s}{1}{x}]$
  ($\Subst{s}{3}{x}$ zou eveneens volstaan). Dus geldt
  $\Sat/{M}{\lnot(R(b,x) \lor R(x,b))}[\Subst{s}{1}{x}]$, en daarmee
  $\Sat/{M}{\lforall[x][\lnot(R(b,x) \lor R(x,b))]}[s]$, en
  bijgevolg $\Sat{M}{\lnot\lforall[x][\lnot(R(b,x) \lor
  R(x,b))]}[s]$. Daarentegen geldt
  \[
  \Sat/{M}{\lexists[x][(R(b,x) \land R(x,b))]}[s].
  \]
  Dat komt doordat $\Sat{M}{\lforall[x][\lnot(R(b,x) \land
  R(x,b))]}[s]$: voor geen enkele $m \in \Domain M$ geldt $\Sat{M}{R(b,x) \land
  R(x,b)}[\Subst{s}{m}{x}]$. Zoals deze voorbeelden al doen vermoeden,
  geldt $\Sat{M}{\lexists[x][!A(x)]}[s]$ dan en slechts dan als
  $\Sat{M}{!A(x)}[\Subst{s}{m}{x}]$ voor ten minste één $m \in \Domain
  M$.}{}

\iftag{defAll}{Om vast te stellen of een universeel gekwantificeerde
  !!{formula}~$\lforall[x][!A(x)]$ wordt vervuld, moeten we nagaan
  of $\Sat{M}{\lnot\lexists[x][\lnot !A(x)]}[s]$. Zo geldt
  \[
  \Sat{M}{\lforall[x][(R(x,a) \lif R(a,x))]}[s],
  \]
  Ten eerste geldt $\Sat{M}{R(x,a) \lif R(a,x)}[\Subst{s}{m}{x}]$ voor alle $m \in
  \Domain M$ ($\Sat{M}{R(a,x)}[\Subst{s}{1}{x}]$ en
  $\Sat/{M}{R(x,a)}[\Subst{s}{m}{x}]$ voor $m = 2$, $3$ of~$4$).
  Er is dus geen $m \in \Domain M$ waarvoor
  $\Sat{M}{\lnot(R(x,a) \lif R(a,x))}[\Subst{s}{m}{x}]$ geldt, en daarom
  $\Sat/{M}{\lexists[x][\lnot(R(x,a) \lif R(a,x))]}[s]$. Bijgevolg
  geldt $\Sat{M}{\lnot\lexists[x][\lnot(R(x,a) \lif R(a,x))]}[s]$. Aan de
  andere kant geldt
  \[
  \Sat/{M}{\lforall[x][(R(a,x) \lif R(x,a))]}[s],
  \]
  want $\Sat/{M}{R(a,x) \lif R(x,a)}[\Subst{s}{2}{x}]$, en dus
  $\Sat{M}{\lexists[x][\lnot(R(a,x) \lif R(x,a))]}[s]$. Zoals deze
  voorbeelden al doen vermoeden, geldt $\Sat{M}{\lforall[x][!A(x)]}[s]$
  dan en slechts dan als $\Sat{M}{!A(x)}[\Subst{s}{m}{x}]$ voor elke $m \in \Domain
  M$.}{}

Beschouw als ingewikkelder geval
\[
\lforall[x][(R(a,x) \lif \lexists[y][R(x,y)])].
\]
Omdat $\Sat/{M}{R(a,x)}[\Subst{s}{3}{x}]$ en
$\Sat/{M}{R(a,x)}[\Subst{s}{4}{x}]$, zijn alleen $m = 1$
en $m = 2$ de relevante gevallen waarin we naar het consequens van de
implicatie moeten kijken. Geldt $\Sat{M}{\lexists[y][R(x,y)]}[\Subst{s}{1}{x}]$?
Dit is zo als er ten minste één $n \in \Domain M$ bestaat waarvoor
$\Sat{M}{R(x,y)}[\Subst{\Subst{s}{1}{x}}{n}{y}]$. Nemen we inderdaad
$n = 1$, dan is $\Subst{\Subst{s}{1}{x}}{n}{y} = \Subst{s}{1}{y} =
s$. Omdat $s(x) = 1$, $s(y) = 1$ en $\tuple{1,1} \in \Assign{R}{M}$,
is het antwoord bevestigend.

Om te bepalen of $\Sat{M}{\lexists[y][R(x,y)]}[\Subst{s}{2}{x}]$ geldt,
moeten we de !!{variable} bedelingen
$\Subst{\Subst{s}{2}{x}}{n}{y}$ bekijken. Voor $n = 1$ is deze bedeling
gelijk aan~$s_2 = \Subst{s}{2}{x}$, die $R(x,y)$ niet vervult ($s_2(x) =
2$, $s_2(y) = 1$ en $\tuple{2,1}\notin \Assign{R}{M}$). Beschouw echter
$\Subst{\Subst{s}{2}{x}}{3}{y} = \Subst{s_2}{3}{y}$.
$\Sat{M}{R(x,y)}[\Subst{s_2}{3}{y}]$, want $\tuple{2,3} \in
\Assign{R}{M}$, en dus $\Sat{M}{\lexists[y][R(x,y)]}[s_2]$.

Voor alle $m \in \Domain M$ geldt dus ofwel
$\Sat/{M}{R(a,x)}[\Subst{s}{m}{x}]$ (als $m = 3$, $4$), ofwel
$\Sat{M}{\lexists[y][R(x,y)]}[\Subst{s}{m}{x}]$ (als $m = 1$, $2$), en daarom
\[
\Sat{M}{\lforall[x][(R(a,x) \lif \lexists[y][R(x,y)])]}[s].
\]
Daarentegen geldt
\[
\Sat/{M}{\lexists[x][(R(a,x) \land \lforall[y][R(x,y)])]}[s].
\]
We hebben $\Sat{M}{R(a,x)}[\Subst{s}{m}{x}]$ alleen voor $m = 1$ en $m =
2$. Voor beide waarden van~$m$ bestaat er op zijn beurt een $n \in
\Domain M$: voor de eerste is dat $n = 4$; voor de tweede is dat het
domeinelement 1. In beide gevallen geldt
$\Sat/{M}{R(x,y)}[\Subst{\Subst{s}{m}{x}}{n}{y}]$ en dus
$\Sat/{M}{\lforall[y][R(x,y)]}[\Subst{s}{m}{x}]$ voor $m = 1$ en $m =
2$. Kortom, er is geen $m \in \Domain M$ waarvoor $\Sat{M}{R(a,x)
\land \lforall[y][R(x,y)]}[\Subst{s}{m}{x}]$.
\end{ex}

\iftag{defEx}{%
\begin{prop}\ollabel{prop:sat-ex}
$\Sat{M}{\lexists[x][!B(x)]}[s]$ dan en slechts dan als er een $x$-variant $s'$ van $s$
    bestaat waarvoor $\Sat{M}{!B(x)}[s']$.
\end{prop}

\begin{proof}
  Oefening.
\end{proof}
}{}

\tagprob{defEx}
\begin{prob}
  Bewijs \olref[fol][syn][sat]{prop:sat-ex}
\end{prob}
\tagendprob

\iftag{defAll}{%
\begin{prop}\ollabel{prop:sat-all}
$\Sat{M}{\lforall[x][!B(x)]}[s]$ dan en slechts dan als voor elke $x$-variant~$s'$ van $s$
  geldt dat $\Sat{M}{!B(x)}[s']$
\end{prop}

\begin{proof}
  Oefening.
\end{proof}
}{}

\tagprob{defAll}
\begin{prob}
  Bewijs \olref[fol][syn][sat]{prop:sat-all}
\end{prob}
\tagendprob

\begin{prob}
Zij $\Lang L = \{c, f, A\}$ met één !!{constant}, één eenplaatsige
!!{function} en één tweeplaatsig !!{predicate}, en zij de
!!{structure}~$\Struct{M}$ gegeven door
\begin{enumerate}
\item $\Domain M = \{1, 2, 3\}$
\item $\Assign{c}{M} = 3$
\item $\Assign{f}{M}(1) = 2, \Assign{f}{M}(2) = 3, \Assign{f}{M}(3) = 2$
\item $\Assign{A}{M} = \{\tuple{1, 2}, \tuple{2, 3}, \tuple{3, 3}\}$
\end{enumerate}
(a) Zij $s(v) = 1$ voor alle !!{variable}s~$v$. Bepaal of
\[
\Sat{M}{\lexists[x][(A(f(z), c) \lif \lforall[y][(A(y, x) \lor A(f(y),
      x))])]}[s]
\]
geldt. Licht uw antwoord toe.

(b) Geef een andere structuur en !!a{variable} bedeling waarin de
!!{formula} niet wordt vervuld.
\end{prob}

\end{document}
